\chapter{Review of Literature}
In this chapter include a critical appraisal of the previous work published in the literature pertaining to the topic of the investigation. This review examines recent advances in machine learning-based URL threat detection and phishing prevention systems.

\section{Review of Paper 1: Efficient Chrome Extension for Phishing Detection Using Machine Learning Techniques} 

\subsection*{Methodology and Approach}
Tha¸ci et al. (2024) developed a Chrome extension called "NoPhish" that employs machine learning techniques for real-time phishing detection. Their approach utilized a dataset from PhishTank containing 22 features rated by Alexa database, implementing three classification algorithms: Random Forest, Support Vector Machine (SVM), and k-Nearest Neighbor (kNN). The system operates as middleware between users and potentially malicious websites, providing immediate threat assessment during web browsing.\cite{arxiv5}

\subsection*{Strengths and Contributions}
\begin{itemize}
	\item Practical implementation as a browser extension for real-world deployment
	\item Comprehensive feature extraction covering URL characteristics, domain properties, and content analysis
	\item Comparative evaluation of multiple machine learning algorithms with Random Forest achieving best precision
	\item Integration of both heuristic analysis and machine learning for enhanced detection capabilities
	\item Focus on user accessibility regardless of cybersecurity expertise level
\end{itemize}

\subsection*{Limitations and Gaps}
\begin{itemize}
	\item Limited to Chrome browser ecosystem, restricting cross-platform compatibility
	\item Dependency on PhishTank dataset may introduce bias toward known attack patterns
	\item No discussion of computational overhead or real-time performance metrics
	\item Lack of evaluation against zero-day phishing attacks or sophisticated evasion techniques
	\item Absence of community feedback mechanisms for continuous model improvement
\end{itemize}

\section{Review of Paper 2: Website Phishing Detection Using Machine Learning Techniques}

\subsection*{Methodology and Approach}
Alazaidah et al. (2024) conducted a comprehensive comparison of 24 different classifiers across six learning strategies for phishing website detection. Their methodology evaluated classifiers including Bayesian methods, function-based approaches, lazy learning, meta-learning, rule-based systems, and tree-based algorithms. The study employed two datasets with different characteristics and assessed performance using eight evaluation metrics including accuracy, precision, recall, F-score, and ROC area.\cite{arxiv}

\subsection*{Strengths and Contributions}
\begin{itemize}
	\item Systematic evaluation of multiple machine learning paradigms providing comprehensive performance comparison
	\item Use of diverse evaluation metrics ensuring robust performance assessment
	\item Analysis of feature selection methods including InfoGainAttributeEval for optimal feature identification
	\item Demonstration of Random Forest, FilteredClassifier, and J-48 superiority in phishing detection
	\item Consideration of both binary and multi-class classification scenarios
\end{itemize}

\subsection*{Limitations and Gaps}
\begin{itemize}
	\item Limited dataset size potentially affecting generalization capabilities
	\item Lack of real-time performance evaluation and computational complexity analysis
	\item No consideration of adversarial attacks or evasion techniques
	\item Absence of hybrid or ensemble approaches beyond individual classifier comparison
	\item Limited discussion of deployment considerations and scalability issues
\end{itemize}

\section{Review of Paper 3: Website Phishing Attack Detection Using Innovative Meta Learning-Based Ensemble Approach}.

\subsection*{Methodology and Approach}
Naseeb et al. (2025) proposed a novel stacked ensemble model combining Artificial Neural Networks (ANN) and Bagging K-Nearest Neighbors (KNN) with Logistic Regression as a meta-learner classifier. Their hybrid framework leveraged ANN's capability for global non-linear pattern identification and Bagging KNN's proficiency in local pattern detection. The system incorporated SMOTE oversampling for class imbalance handling and employed k-fold cross-validation for robust performance validation using 2019-2024 datasets.\cite{IEEE}

\subsection*{Strengths and Contributions}
\begin{itemize}
	\item Innovative stacking ensemble approach combining complementary machine learning paradigms
	\item Impressive performance metrics achieving 97.20\% accuracy, 97\% precision, and 97\% recall
	\item Comprehensive preprocessing pipeline including SMOTE for class imbalance mitigation
	\item Real-time optimization using TensorFlow Lite for edge-level deployment feasibility
	\item Use of contemporary datasets spanning 2019-2024 ensuring relevance to current threat landscape
	\item Detailed hyperparameter optimization with 512-256-128-128 layer ANN architecture
\end{itemize}

\subsection*{Limitations and Gaps}
\begin{itemize}
	\item High computational complexity may limit real-time deployment in resource-constrained environments
	\item Dependency on balanced datasets through SMOTE may not reflect real-world class distributions
	\item Limited evaluation against adversarial examples or sophisticated attack variations
	\item Lack of comparison with recent deep learning approaches such as transformer-based models
	\item No discussion of model interpretability or explainability for security practitioners
\end{itemize}
